\documentclass[10pt,a4paper]{amsart}
\usepackage[margin=19mm]{geometry}
\usepackage{amsmath,amssymb,mathtools}
\usepackage[scr=rsfs,cal=euler]{mathalfa}
\usepackage{xfrac}
\usepackage[unicode,hidelinks,bookmarks=false]{hyperref}
\newcommand{\ie}{i.e.\ }
\newcommand{\cf}{cf.\ }
\newcommand{\resp}{respectively\ }
\newcommand{\Subsection}{\subsection}
\providecommand{\nobreakdash}{\nobreak\hbox{-}}
\setlength{\emergencystretch}{2em}
\multlinegap0pt
\usepackage{bm}
\usepackage{bbm}
\usepackage{dsfont}
\usepackage{xspace}
\usepackage{cancel}
\usepackage{mathtools}
\usepackage{amsthm}
\makeatletter
\@ifundefined{theorem}{\newtheorem{theorem}{Theorem}}{}
\@ifundefined{claim}{\newtheorem{claim}{Claim}}{}
\@ifundefined{lemma}{\newtheorem{lemma}[theorem]{Lemma}}{}
\makeatother
\title[Two results on set families: sturdiness and intersection]{Two results on set families: sturdiness and intersection}
\author{Yongjiang Wu, Zhiyi Liu, Lihua Feng, Yongtao Li}
\date{Source: arXiv:2508.19723v2 (CC BY 4.0)}
\begin{document}
\maketitle

\noindent\emph{Source and licence.} Two results on set families: sturdiness and intersection. Version arXiv:2508.19723v2 (CC BY 4.0), distributed under the Creative Commons Attribution 4.0 International licence, as recorded in the arXiv licence metadata for this exact version and shown on the abstract page https://arxiv.org/abs/2508.19723v2. Full rights notice: licenses/arxiv-2508.19723-CC-BY-4.0.txt.\par\smallskip

\noindent\textit{Editorial scope.} Counted unit: the Lemma whose source label is \texttt{35} in the article (source0.tex lines 486--494, source label \texttt{35}), delivered together with the article's own complete original proof of it (source0.tex lines 495--594, one \texttt{proof} environment of 100 lines). No mathematics is written, repaired or re-derived: statement and proof are the article's own text line for line. Required context retained uncounted: 34 (lines 429--435). Every definition, display and label of those lines is retained unchanged. Omitted from this package and not redistributed: 1--428, 436--485, 595--880. Nothing inside a retained excerpt is omitted or altered: every retained body is delivered exactly as the article's lines are written, so no omission receipt of any kind accompanies this package. Citations: the retained excerpts carry no citation command at all, so no bibliography entry of the article is transcribed and no reference is redistributed. Reference adaptation: none was needed and none was made; every reference inside the delivered unit resolves inside it, so no unresolved reference or citation is produced. Printed slips kept literal: no printed word, symbol, hypothesis or number of the counted statement or of its proof was altered. Layout and class: the article's own class is \texttt{article} with packages including amsmath, amsfonts, amsthm, amssymb, mathtools, mathrsfs, graphicx, color, latexsym, tikz, calc, tikz-cd, epstopdf, enumerate; the wrapper is the lane's pinned \texttt{amsart} editorial wrapper at 10pt on A4, which restates the theorem-like environments the retained text needs and adds the article's own macro definitions that the retained text needs (none), with extra wrapper packages (bm, bbm, dsfont, xspace, cancel, mathtools). The delivered numbering is the unit's own. Assets: no retained span contains a figure environment, a graphics-inclusion command, a PDF-page inclusion, a table or a TikZ picture; no image file is copied into this package, the package declares no asset dependency and no image file is redistributed with it. Licence: version arXiv:2508.19723v2 of this article is distributed under the Creative Commons Attribution 4.0 International licence, as recorded in the arXiv licence metadata for this exact version and shown on the abstract page https://arxiv.org/abs/2508.19723v2. A full rights notice is delivered at licenses/arxiv-2508.19723-CC-BY-4.0.txt. Characters: every retained excerpt is pure ASCII, so the delivered engine, XeLaTeX with its default Latin Modern text font, reports no missing character.
 \begin{lemma}\label{34}
      Let $\mathcal{F}, \mathcal{G} \subseteq 2^{[k]}$ be cross $t$-intersecting families and $1\le i< j\le k$. Suppose that $a$ is the maximal necessary intersection point of $\mathcal{F}$ and $ \mathcal{G}$. Then 
      $
       \left|[a-1] \cap F\cap G\right|\ge t-1.
      $
 Moreover, $s_{ij}(\mathcal{F})$ and $ s_{ij}(\mathcal{G})$ are  cross $t$-intersecting with the maximal necessary intersection point at most $a$. 
 \end{lemma}

\begin{lemma}\label{35}
     Let $k \geq \ell \geq \ell' \geq t\geq 1$  be integers. For $i \in[2]$, let $\omega_i:[0,k] \rightarrow \mathbb{R}_{\geq 0}$ be non-increasing functions.  Let $\mathcal{F}\subseteq\binom{[k]}{\leq \ell}$ and $\mathcal{G} \subseteq\binom{[k]}{\leq \ell'}$ be non-empty cross $t$-intersecting families with the maximal necessary intersection point $a$. Suppose that $a \geq t+1$ and $\mathcal{F} \backslash \mathcal{F}^a \neq \emptyset$ and $\mathcal{G} \backslash \mathcal{G}^a \neq \emptyset$. Then there are families $\mathcal{F}^*\subseteq\binom{[k]}{\leq \ell}$ and $\mathcal{G}^* \subseteq\binom{[k]}{\leq \ell'}$  such that
     \begin{itemize}
         \item [(1)] $\mathcal{F}^*, \mathcal{G}^*$ are non-empty cross $t$-intersecting families;
        \item[(2)]  $\omega_1(\mathcal{F}^*)+\omega_2(\mathcal{G}^*)\ge \omega_1(\mathcal{F})+\omega_2(\mathcal{G})$;
        \item[(3)] the maximal necessary intersection point of $\mathcal{F}^*$ and $\mathcal{G}^*$ is smaller than $a$.
         \item[(4)] there are shifted families $\mathcal{F}^*$ and $\mathcal{G}^*$ satisfying (1)-(3).
     \end{itemize}
\end{lemma}

\begin{proof}
Observe that (1) and (2) are preserved under shifting. According to Lemma \ref{34}, 
(3) is also preserved under shifting.
 Since one can apply the shifting operator to make $\mathcal{F}^*\subseteq\binom{[k]}{\leq \ell}$ and $\mathcal{G}^* \subseteq\binom{[k]}{\leq \ell'}$ shifted families, it suffices to establish properties (1)-(3).

Define
    $
    \mathcal{F}^{add}=\left\{F\backslash\{a\} : F\in \mathcal{F}^a   \right\},\ \mathcal{G}^{add}=\left\{G\backslash\{a\} : G\in \mathcal{G}^a   \right\}.
    $
Furthermore, let 
$$\mathcal{F}^+=\mathcal{F}\cup\mathcal{F}^{add},\quad \mathcal{G}^+=\mathcal{G}\cup\mathcal{G}^{add},\quad \mathcal{F}^-=\mathcal{F}\backslash\mathcal{F}^{a},\quad \mathcal{G}^-=\mathcal{G}\backslash\mathcal{G}^{a}.$$
Then  $\mathcal{F}^+, \mathcal{F}^-\subseteq \binom{[k]}{\leq \ell}$ and 
$\mathcal{G}^+, \mathcal{G}^-\subseteq \binom{[k]}{\leq \ell'}$.
\begin{claim}\label{cl1}
$\mathcal{F}\cap\mathcal{F}^{add}=\emptyset$ and $\mathcal{G}\cap\mathcal{G}^{add}=\emptyset$.
    \end{claim}
    \begin{proof}[Proof of Claim \ref{cl1}]
By symmetry, it suffices to prove that $\mathcal{F}\cap\mathcal{F}^{add}=\emptyset$.
        Assume, for the sake of contradiction, that  there exists $F\in\mathcal{F}\cap\mathcal{F}^{add}$.
     Since $F\in\mathcal{F}^{add}$, there is $F^*\in\mathcal{F}^a$
such that $F=F^*\backslash \{a\}$. By the definition of 
$\mathcal{F}^a$, there exists  $G\in \mathcal{G}$
 such that 
        $|[a]\cap F^*\cap G|=t$
        and $a\in F^*\cap G$.
This implies that $|[a-1]\cap F^*\cap G|=t-1$ and hence
        $
        |[a-1]\cap F\cap G|=t-1.
        $
Since $a\notin F$,  we obtain $|[a]\cap F\cap G|=t-1$, which contradicts with the maximality of $a$.
    \end{proof}

Since $ \omega_1, \omega_2$ are non-increasing, we obtain $\omega_1(\mathcal{F}^{add})\geq \omega_1(\mathcal{F}^a)$ and $\omega_2(\mathcal{G}^{add})\geq \omega_2(\mathcal{G}^a)$.
If $\omega_1(\mathcal{F}^a)\geq \omega_2(\mathcal{G}^a)$, then let
 $\mathcal{F}^*= \mathcal{F}^+$ and $\mathcal{G}^*=\mathcal{G}^-$.
It follows from Claim \ref{cl1} that
 \begin{align*}
       \omega_1(\mathcal{F}^*)+\omega_2(\mathcal{G}^*)&= \omega_1(\mathcal{F})+\omega_2(\mathcal{G})+\omega_1(\mathcal{F}^{add})-\omega_2(\mathcal{G}^a)\\
       &\geq\omega_1(\mathcal{F})+\omega_2(\mathcal{G})+\omega_1(\mathcal{F}^{a})-\omega_2(\mathcal{G}^a)\\
&\geq \omega_1(\mathcal{F})+\omega_2(\mathcal{G}).
        \end{align*}
If $\omega_1(\mathcal{F}^a)\leq \omega_2(\mathcal{G}^a)$, then let
 $\mathcal{F}^*= \mathcal{F}^-$ and $\mathcal{G}^*=\mathcal{G}^+$.
Similarly, we have $\omega_1(\mathcal{F}^*)+\omega_2(\mathcal{G}^*)\geq \omega_1(\mathcal{F})+\omega_2(\mathcal{G})$. Thus (2) holds.

It remains  to prove the following claim. 
 \begin{claim}\label{cl2}
  For any $F^+\in\mathcal{F}^+$ and $G^-\in\mathcal{G}^-$, 
     $
    |[a-1] \cap F^+\cap G^-|\geq t
     $
     holds.
Moreover, for any $F^-\in\mathcal{F}^-$ and $G^+\in\mathcal{G}^+$, 
     $
    |[a-1] \cap F^-\cap G^+|\geq t
     $
     holds.
    \end{claim}
    \begin{proof}[Proof of Claim \ref{cl2}]
By symmetry, it suffices to prove the first assertion.
    Suppose that there exist $F^+\in\mathcal{F}^+$ and $G^-\in\mathcal{G}^-$ such that
     $
    |[a-1] \cap F^+\cap G^-|< t.
     $


  If $F^+\in\mathcal{F}$,  then $G^-\in\mathcal{G}$ and
Lemma \ref{34} imply $
    |[a-1] \cap F^+\cap G^-|=t-1.
     $
By the maximality of $a$, we have $
    |[a] \cap  F^+\cap G^-|=t
     $
and $a\in F^+\cap G^-$.
Then $ G^-\in \mathcal{G}^a$,  a contradiction.

If $F^+\notin\mathcal{F}$, then $F^+\in\mathcal{F}^{add}$.  Since  $F^+\in\mathcal{F}^{add}$, there exists 
 $F\in\mathcal{F}^a$ such that
        $ F^+=F\backslash \{a\}$.
In view of    $
          |[a-1] \cap  F^+\cap G^-|< t,
        $
it follows that
      $
        \left|[a-1]\cap F\cap G^- \right|<t.
       $
Then $G^-\in\mathcal{G}$ and
Lemma \ref{34} imply 
$
        \left|[a-1]\cap F\cap G^- \right|=t-1.
       $
By the maximality of $a$, we have $
    |[a] \cap F\cap G^-|=t
     $
and $a\in  F\cap G^-$,
       which contradicts with  $ G^-\in \mathcal{G}\backslash\mathcal{G}^a$.
    \end{proof}

Since  Claim \ref{cl2} implies that (1) and (3) hold, the proof of the lemma is finished. 
\end{proof}

\end{document}
